\section{Datasets}

% \begin{table}[h]
% \centering
% \caption{Dataset characteristics.}
% \label{tab:datasets}
% \setlength{\tabcolsep}{4pt} % Tights column spacing to fit width
% \begin{tabular}{l|cccc}
% \toprule
% \textbf{Dataset} & \textbf{Patients} & \textbf{Scans} & \textbf{Pathology} & \textbf{Region}\\
% \midrule
% autoPET IV & 285 & 670 & Melanoma Lesions & Whole-Body \\
% PanTrack (Ours) & 45 & 161 & Pancreatic Cancer & Abdomen \\
% \bottomrule
% \end{tabular}
% \end{table}

\noindent \textbf{autoPET/CT IV Dataset.} This dataset comprises longitudinal whole-body CT scans from 285 melanoma patients undergoing therapy response assessment~\cite{kuestner2025longitudinalct} totalling 670 images. Each patient has at least one baseline and follow-up scan acquired during portal-venous phase on multiple Siemens scanners with standardized protocols. Two radiologists manually segmented all tumor lesions across timepoints with side-by-side verification to establish correspondence. The dataset includes challenging scenarios (splitting, merging, vanishing lesions) and provides pre-computed lesion centers at both timepoints, with registration-based propagated clicks simulating the \textit{Verified Tracking} workflow.\\

\noindent \noindent \textbf{The PanTrack Dataset.} To evaluate generalization to a different anatomical domain, we curated \textit{PanTrack}, comprising 45 patients with pancreatic adenocarcinoma amounting to 161 CT examinations (2–11 per patient; mean 3.6). All scans were acquired on identical Siemens protocols (portal-venous phase) at a single institution. An experienced radiologist with pancreatic imaging expertise manually segmented all pancreatic lesions across timepoints, including hepatic metastases if present. The cohort represents diverse trajectories: some patients underwent long-term stable chemotherapy, others showed rapid progression. Unlike melanoma metastases, pancreatic lesions exhibit fuzzy boundaries and subtle soft-tissue contrast, providing complementary evaluation under different radiological characteristics. This dataset is publicly released upon acceptance.